\documentclass[a4paper,11pt]{ctexbook}
\usepackage{finance-style}

\title{\textcolor{themeblue}{\Huge\textbf{金融学基础}}}
\subtitle{\textcolor{gray}{第一卷\quad 共八课}}
\author{}
\date{}

\begin{document}
\maketitle
\thispagestyle{empty}
\tableofcontents
\newpage

% ============================================================
\chapter{货币的时间价值——一切定价的起点}
% ============================================================

\begin{scenario}
\textbf{北京家庭的教育金规划}

张阿姨的儿子今年8岁。她计划10年后送孩子上大学，届时需要一笔约50万元的教育金。目前她有20万元存款可一次性投入，同时每月还能结余3000元用于储蓄。她面临的选择：存银行定期（年利率1.5\%）、买国债（2.5\%）、买银行理财（3\%）、或配置一部分股票指数基金（预期收益8\%但波动大）。哪个方案能达到目标？差多少？这引出了金融学最基本的问题：\textbf{今天的钱和未来的钱如何比较？}
\end{scenario}

\begin{qualitative}
\textbf{为什么今天的100元比明天的100元更值钱？}

原因有三：\textbf{机会成本}——今天的钱可以投资产生回报；\textbf{通胀}——物价上涨导致未来同样金额的购买力下降；\textbf{不确定性}——未来的承诺未必能兑现。这三个因素共同决定了"货币的时间价值"（Time Value of Money, TVM）——它是整个金融定价体系的基石。无论是债券、股票、衍生品还是房地产，所有资产的定价都归结为同一件事：把未来的现金流折现到今天。
\end{qualitative}

\begin{quantitative}
\textbf{终值与现值}

将$PV$元以年利率$r$投资$n$期，按年复利：
\[
FV = PV \times (1 + r)^n
\]

反过来，未来$n$期的$FV$元在今天的价值：
\[
PV = \frac{FV}{(1 + r)^n}
\]

折现率$r$每增加一个百分点，现值呈指数级下降。这就是"折现"的含义——未来的钱要打折扣才能换算到今天。

\textbf{实际年利率（EAR）与名义年利率（APR）：}如果一年复利$m$次：
\[
EAR = \left(1 + \frac{APR}{m}\right)^m - 1
\]

当$m \to \infty$时得到连续复利：$FV = PV \cdot e^{rt}$。

\textbf{年金公式}——等额定期现金流的终值：
\[
FV = PMT \times \frac{(1 + r)^n - 1}{r}
\]

现值：
\[
PV = PMT \times \frac{1 - (1 + r)^{-n}}{r}
\]

\textbf{72法则}——近似估算翻倍时间：
\[
t \approx \frac{72}{r \times 100}
\]
当$r=8\%$时，约9年翻倍。精确解：$\ln2 / \ln1.08 \approx 8.99$年。

\textbf{案例计算：} 张阿姨20万元一次性投资10年。\\
若存银行（1.5\%）：$200000 \times 1.015^{10} \approx 232,100$元——远不够50万。\\
若年化6\%（合理的股债组合）：$200000 \times 1.06^{10} \approx 358,200$元——仍不够。\\
若每月定投3000元年化6\%：$3000 \times \frac{(1.005)^{120} - 1}{0.005} \approx 491,000$元——加上本金355,200+491,000≈846,200元。远超目标。\\
\textbf{结论：}一次性投入+长期定投，在合理收益率下完全可达目标。
\end{quantitative}

\begin{lessonsummary}
\begin{enumerate}
  \item 货币的时间价值是金融定价的基石——一切资产的价值等于其未来现金流的现值之和
  \item 复利让时间成为投资者最强大的工具，也是最危险的敌人（对债务人而言）
  \item 72法则提供了快速估算的直觉工具，在5-15\%收益率范围内精确度很高
  \item 年金公式是退休规划和贷款摊销的核心工具
\end{enumerate}
\end{lessonsummary}

\begin{discuss}
\begin{enumerate}
  \item "72法则在低利率时误差大"——当$r=1\%$时，72法则给出72年，精确值为69.7年。为什么误差随着利率降低而增大？
  \item 折现率的选取本质上是主观的。两个人对同一笔未来现金流的估值可能相差30\%以上——你认为"一个人的主观折现率"受哪些因素影响？（提示：耐心、财富水平、生命阶段、文化背景等）
  \item 负利率（如2010年代欧洲和日本）意味着什么？为什么有人愿意购买负利率债券——明知道到期会亏钱？
\end{enumerate}
\end{discuss}

\begin{exercises}
\begin{enumerate}
  \item 某理财产品宣称年化收益率5\%，但按季度复利。它的EAR是多少？如果按月复利呢？
  \item 你25岁，希望在60岁退休时有500万。假设年化收益7\%，你每月需要定投多少？如果从35岁才开始呢？差了10年，每月需要多投多少？（提示：用年金现值公式）
  \item 两个互斥项目：A项目投资100万，未来3年现金流分别为50万、50万、50万；B项目投资100万，现金流分别为0、0、200万。折现率10\%，分别计算NPV和IRR。哪个项目更优？如果折现率变成5\%或15\%呢？结论对折现率敏感吗？
\end{enumerate}
\end{exercises}

\xuehaiinline{
\textbf{贴现率的哲学争论——Ramsey方程与社会贴现率：}经济学家在评估跨代项目（如气候变化政策）时面临一个伦理难题——应该用多高的折现率？高折现率意味着"未来人的福祉不如当代人重要"，低折现率则相反。Ramsey（1928）提出$r = \delta + \eta g$，其中$\delta$是纯时间偏好率、$\eta$是边际效用弹性、$g$是人均消费增长率。如果用$\delta=0$（当代人和未来人同样重要），社会贴现率大约在1-2\%——这远低于金融市场上的私人贴现率。Stern Review（2006）用低贴现率得出了"立即大力减排"的结论；Nordhaus（反对者）用市场贴现率得出了"逐步减排"的结论。同一个问题，不同的贴现率假设，决定了完全不同的政策方向。\\
\textbf{72法则的"精确版"：}严格推导——$(1+r)^t = 2 \implies t = \ln2 / \ln(1+r)$。当$r$很小时，$\ln(1+r) \approx r - r^2/2$。代入：$t \approx \ln2 / (r - r^2/2) = (0.693 / r) \times 1/(1 - r/2)$。这就是为什么69.3比72在小利率时更精确——但72在常见利率范围（5-15\%）内更易心算。
}

\newpage
% ============================================================
\chapter{债券定价与利率风险}
% ============================================================

\begin{scenario}
\textbf{最安全资产的不安全时刻}

2020年，一位投资者以面值100元买入10年期中国国债，票面利率3\%。到2023年，央行降息后新发行的同类国债票面利率降至2.5\%。他的债券现在值多少钱？——超过100元，因为老债券每年多付0.5元利息。但如果央行加息呢？老债券就会贬值。2022年至2023年，美联储加息500bp，长期美国国债价格暴跌——持有大量美国国债的硅谷银行（SVB）因此浮亏约180亿美元，最终在2023年3月倒闭。持有"最安全"的资产，却死于"最安全"的利率风险。
\end{scenario}

\begin{qualitative}
\textbf{债券定价的核心逻辑}

债券的本质是一张标准化借条。发行人承诺在特定日期支付固定利息（票息）并偿还本金。其价格就是未来所有现金流按市场利率折现的现值之和。

\textbf{折价、平价与溢价：}当市场利率 > 票面利率时，债券以低于面值的价格交易（折价）。因为投资者可以买到利率更高的新债券，老债券必须降价来补偿。反之则溢价交易。

\textbf{最关键的关系：利率和债券价格反向变动。}这个反向关系不是线性的——久期（Duration）衡量的是价格对利率变化的敏感度。
\end{qualitative}

\begin{quantitative}
\textbf{债券定价公式}

面值$F$、票面利率$c$、每年付息$C = cF$、$n$年期、市场到期收益率$y$：
\[
P = \sum_{t=1}^{n} \frac{C}{(1+y)^t} + \frac{F}{(1+y)^n}
\]

\textbf{到期收益率（YTM）：}使债券当前市价等于未来现金流现值之和的折现率——是债券"真实"年化回报的衡量。

\textbf{久期（Macaulay Duration）：}
\[
D = \frac{\sum_{t=1}^{n} t \times \frac{C}{(1+y)^t} + n \times \frac{F}{(1+y)^n}}{P}
\]

\textbf{修正久期（Modified Duration）：}$D^* = D / (1+y)$。价格变化近似：
\[
\Delta P / P \approx -D^* \times \Delta y
\]

\textbf{凸性（Convexity）：}久期是线性近似，凸性校正二阶项：
\[
\Delta P / P \approx -D^* \times \Delta y + \frac{1}{2} \times Convexity \times (\Delta y)^2
\]

久期越长的债券对利率越敏感。零息债券的久期等于其剩余期限——是普通付息债券的1.3-1.8倍。
\end{quantitative}

\begin{lessonsummary}
\begin{enumerate}
  \item 债券价格 = 未来票息和本金的现值之和，与市场利率反向变动
  \item 久期衡量债券价格对利率的敏感度——是风险管理的核心工具
  \item 凸性校正了久期的线性近似误差，在利率大幅变动时必须考虑
  \item 利率风险不仅影响高风险债券——它同样作用于"最安全"的国债
\end{enumerate}
\end{lessonsummary}

\begin{discuss}
\begin{enumerate}
  \item 硅谷银行（SVB）持有大量"持有至到期"（HTM）的美国国债，按会计准则不需要按市值计价。这是否帮助管理层"隐藏"了风险？会计准则应当在银行风险管理中扮演什么角色？
  \item 2023年SVB倒闭后，美国监管机构推出了银行期限资金计划（BTFP）——允许银行按面值而不是按市值向美联储质押国债借款。这个政策是在解决什么问题？它是否有道德风险？
  \item 如果一家银行的资产久期是5年、负债久期是1年，利率上升200bp时净资产的变化是多少？中国的商业银行普遍存在"期限错配"（借短贷长）——为什么它们没有像SVB一样倒闭？
\end{enumerate}
\end{discuss}

\begin{exercises}
\begin{enumerate}
  \item 一张面值100元、票面利率4\%、3年期的债券，每年付息。当前市场YTM为3\%和6\%时，价格分别是多少？（精确到分）
  \item 计算上述债券在YTM=4\%时的Macaulay久期和修正久期。如果YTM从4\%上升到5\%，用修正久期预估价格变化，再精确计算实际价格变化。二阶近似的凸性校正需要加多少？
  \item 两个债券：A为10年期国债（票面3\%），B为10年期零息国债。分别计算它们的久期。如果市场利率上升100bp，哪个跌得多？为什么？
\end{enumerate}
\end{exercises}

\xuehaiinline{
\textbf{硅谷银行倒闭的完整技术链：（1）SVB负债端——创业公司存款（2020-2021年科技繁荣期大量流入），其中94\%超过FDIC保险上限（25万美元）。这些"大额未保险存款"对利率极其敏感；（2）资产端——SVB将大部分存款投入了长期（5-10年）美国国债和MBS，分类为HTM，账面价值不变；（3）美联储在2022年加息500bp——SVB的HTM资产市值蒸发约180亿美元，但资产负债表上没有体现；（4）2023年3月SVB宣布需要融资补充资本金——触发了存款人的恐慌。通过网上银行，存款人在几小时内取走420亿美元（25\%的存款）；（5）SVB被迫出售HTM资产——按市值变现损失约18亿美元——资本金不足以覆盖，48小时内被FDIC接管。\\
\textbf{教训：}久期缺口管理不当 + 存款集中度过高 + HTM会计分类是"三条腿的凳子"——任何一条断裂都会导致崩塌。\\
\textbf{中国的商业银行体系：}大型银行（工农中建）存款来源高度分散且稳定（零售存款占比>50\%），贷款以浮动利率重定价为主（久期短），加上央行流动性支持工具充足——这些因素使中国银行体系对利率冲击的耐受力远强于SVB。
}

\newpage
% ============================================================
\chapter{收益率曲线与期限结构}
% ============================================================

\begin{scenario}
\textbf{收益率倒挂——最准确的经济衰退预警信号}

2024年8月，中国10年期国债收益率2.1\%，2年期1.5\%——长端高于短端，这是正常形态。但在2022年，美国10年期与2年期国债收益率发生了"倒挂"（10年收益率低于2年）——此后所有期限都出现倒挂。历史数据显示，过去60年里美国每次经济衰退之前都出现了收益率曲线倒挂，平均领先12-18个月。为什么"长端利率低于短端"这个反常现象能如此准确地预测衰退？
\end{scenario}

\begin{qualitative}
\textbf{为什么长期利率一般高于短期利率？}

三个主要理论解释收益率曲线的形状：

\textbf{流动性偏好理论（Liquidity Preference Theory）：}长期债券的风险更高（久期更长、通胀和利率不确定性更大），投资者要求"期限溢价"作为补偿。

\textbf{预期理论（Expectations Theory）：}长期利率等于未来短期利率预期的平均值。收益率曲线向上倾斜意味着市场预期未来短期利率会上升。

\textbf{市场分割理论（Segmented Markets Theory）：}不同期限的债券有不同的主要投资者（养老金/保险偏好长期、银行/货币基金偏好短期），供需决定了各期限的利率。

当代金融学认可的是三者融合的\textbf{期限结构综合理论}。
\end{qualitative}

\begin{quantitative}
\textbf{即期利率与远期利率}

即期利率$r_t$：当前市场上$t$年期零息债券的到期收益率。

远期利率$f_{t,T}$：从未来$t$时刻到$T$时刻的利率，由即期利率隐含：
\[
(1 + f_{t,T})^{T-t} = \frac{(1 + r_T)^T}{(1 + r_t)^t}
\]

远期利率可以理解为"市场对未来利率的隐含预期"加上一定的期限溢价。

\textbf{Bootstrapping——从附息债券推导即期利率曲线：}

给定一组不同期限的附息国债价格，通过逐次求解获得各期限的即期利率。这是收益率曲线构建的基础方法。

\textbf{利率期限溢价的估计：}期限溢价 = 远期利率 - 预期未来即期利率。实证研究（如Kim-Wright, 2005）表明期限溢价随时间显著变化，可从0到200bp不等。
\end{quantitative}

\begin{lessonsummary}
\begin{enumerate}
  \item 收益率曲线形态反映了市场对未来利率、通胀和风险的集体判断
  \item 收益率曲线倒挂是经济衰退最可靠的领先指标——因为它直接反映了市场对未来的悲观预期
  \item 远期利率曲线是债券定价、利率互换定价和央行政策评估的核心工具
  \item 中国的收益率曲线和美国的收益率曲线在形成机制上有系统差异
\end{enumerate}
\end{lessonsummary}

\begin{discuss}
\begin{enumerate}
  \item 收益率曲线倒挂为什么能预测衰退？推演逻辑链——市场预期未来经济放缓→预期央行降息→长期利率下降→长端低于短端。但这更像"衰退预期导致倒挂"而非"倒挂导致衰退"——两者方向不要搞反。
  \item 2022年倒挂后美国经济在2023-2024年并没有立即陷入衰退——"软着陆"（通胀下降+就业稳定）实现了。"这次不一样"吗？还是衰退只是被延迟了？
  \item 中国自2010年代后收益率曲线几乎没有倒挂过——是因为中国没有衰退风险，还是因为国债市场的投资者结构（银行持有至到期为主、交易盘占比低）抑制了曲线形态的多样化？
\end{enumerate}
\end{discuss}

\begin{exercises}
\begin{enumerate}
  \item 1年期即期利率2\%，2年期即期利率3\%，计算1年后的1年期远期利率$f_{1,2}$。如果市场对期限溢价估计为0.5\%（1年期），那么市场对未来1年期利率的预期是多少？
  \item 给定三只零息国债：1年期价格98.04（YTM=2\%）、2年期价格94.26（YTM=3\%）、3年期价格90.06（YTM=3.5\%），用bootstrapping推导即期利率曲线。计算$f_{1,2}$和$f_{2,3}$。
  \item 用Excel或Python从央行网站下载近5年某期限利差（10y-2y）的时间序列数据。标注利差倒挂的时期，与同期的GDP增速和股市表现对比。数据来源：中国人民银行官网或Wind。
\end{enumerate}
\end{exercises}

\xuehaiinline{
\textbf{中国的收益率曲线为什么与众不同？}国债市场的主要投资者是商业银行（占60\%+），而银行的持有动机不是交易而是资产负债管理和流动性管理——它们持有至到期，不频繁交易。这导致：（1）中国国债的交易量/存量比远低于美国；（2）收益率曲线的信息含量不如美国——它更多地反映了银行体系的资金状况而非宏观预期；（3）长期国债的期限溢价被低估。2020年央行引入"利率走廊"机制后，短期利率的锚定效应增强，但长端的定价仍然主要受制于供需而非预期。\\
\textbf{Probit模型与衰退预测：}美联储经济学家用Probit模型检验了收益率曲线倒挂对衰退的预测能力——以10年-3个月利差为解释变量。结论：利差转负后的4个季度内出现衰退的条件概率超过80\%。这是宏观金融学中最可靠的实证关系之一。但注意：它预测的是概率，不是确定性——2023年倒挂后没有立即衰退，正说明了"概率"和"必然"之间的差距。
}

\newpage
% ============================================================
\chapter{股票估值——从Gordon模型到DCF}
% ============================================================

\begin{scenario}
\textbf{茅台的估值之谜}

2021年初，贵州茅台的市盈率超过60倍。按当时利润计算，买入后需要60年才能靠利润回本——但投资者仍然蜂拥买入，因为他们相信茅台能保持15\%以上的年增长率。到2024年，茅台市盈率降至约30倍。发生了什么？——市场对茅台未来增长的预期从15\%下调至约8-10\%。增长率预期下降不到一半，市盈率却腰斩。这背后的数学原理就是"增长率的乘法器效应"。
\end{scenario}

\begin{qualitative}
\textbf{股票价值的本质}

股票代表一家公司的所有权。股票的价值等于其未来所有净现金流的现值之和。增长是估值的"放大器"——增长率每变化一个百分点，估值可以变化10-30\%。

股票估值的三种主要方法：
\begin{enumerate}
  \item \textbf{股息贴现（DDM）：}假设价值来自未来股息
  \item \textbf{市盈率法（P/E）：}通过可比公司和行业标准快速判断
  \item \textbf{现金流贴现（DCF）：}最完整的框架——考虑所有自由现金流
\end{enumerate}
\end{qualitative}

\begin{quantitative}
\textbf{Gordon增长模型（DDM）}

假设股息以固定增长率$g$永续增长：
\[
P_0 = \frac{D_1}{r - g} = \frac{D_0(1+g)}{r - g}
\]

从Gordon模型可以推导出市盈率的理论表达式。假设利润留存率为$b$，股本回报率为ROE，则$g = b \times ROE$，$D_1 = EPS_1 \times (1-b)$：
\[
\frac{P_0}{EPS_1} = \frac{1-b}{r - b \times ROE}
\]

\textbf{多阶段增长模型：}高速增长阶段 + 过渡阶段 + 永续增长阶段。

\textbf{DCF模型的企业价值：}
\[
EV = \sum_{t=1}^{n} \frac{FCF_t}{(1+WACC)^t} + \frac{TV}{(1+WACC)^n}
\]

\textbf{终值计算：}永续增长法 $TV = FCF_n(1+g) / (WACC - g)$，或退出倍数法。

\textbf{WACC（加权平均资本成本）：}
\[
WACC = \frac{E}{E+D} \times r_E + \frac{D}{E+D} \times r_D \times (1 - T_c)
\]

\textbf{敏感性分析：}WACC每变动0.5\%，终值增长率每变动0.25\%，DCF估值可以相差20-30\%——这就是为什么不同分析师对同一公司的目标价可以差50\%以上。
\end{quantitative}

\begin{lessonsummary}
\begin{enumerate}
  \item 股票估值本质是对未来现金流的预测和折现——不同的预测产生不同的估值
  \item DCF框架强迫分析者明确做出假设，但假设的微小变化会导致结果的巨大差异——这是精确的错误，不是模糊的正确
  \item 市盈率是市场最常用的快速估值工具，但需结合行业平均和公司增长来判断"高估还是低估"
  \item DCF对高增长、不盈利、轻资产公司的估值几乎完全失效——这些公司的价值在于期权而非现金流
\end{enumerate}
\end{lessonsummary}

\begin{discuss}
\begin{enumerate}
  \item 亚马逊连续亏损十几年，如果1999年用DCF估值几乎必然给出"卖出"——但持有到今天的回报率超过了1000倍。DCF框架的问题出在哪里？"实物期权"（Real Options）能否弥补这个缺陷？
  \item 中国A股市场适用多少倍P/E是"合理"的？一个常见说法是"A股平均市盈率15-20倍是合理的"——但这个"合理"的基准是从哪里来的？美国和日本的历史平均P/E分别是多少？不同国家的P/E为什么有系统性差异？
  \item 茅台从60倍P/E降到30倍——是"估值回归合理"还是"被市场错杀"？试图判断"合理估值"本身就是在做"市场择时"（Market Timing）——而学术研究一致表明市场择时几乎不可能持续成功。
\end{enumerate}
\end{discuss}

\begin{exercises}
\begin{enumerate}
  \item 假设茅台2024年EPS=60元，预期ROE=30\%，利润留存率60\%，必要回报率9\%，用Gordon模型推导合理P/E。如果市场预期ROE降到20\%呢？P/E从多少降到多少？
  \item 给定一家公司的三年财务预测：第1年FCF=1亿，第2年=1.2亿，第3年=1.5亿，终值增长率为3\%，WACC=10\%。计算企业价值。做敏感性分析——WACC从8\%到12\%（步长1\%）、终值g从2\%到4\%（步长0.5\%），给出企业价值的区间。
  \item 下载A股银行板块（如工商银行、招商银行、兴业银行）近5年的P/B和P/E数据，与同期的ROE和不良贷款率对比。为什么银行的估值常用P/B而消费股常用P/E？
\end{enumerate}
\end{exercises}

\xuehaiinline{
\textbf{DCF的七宗罪——为什么"精确"的估值经常错得离谱：}（1）预测偏差——管理层和投资者都有系统性过度乐观的倾向；（2）终值占比过高——在很多DCF中，终值占企业价值的60-80\%，而终值的计算极度依赖一个假设（永续增长率）；（3）WACC的"循环论证"——WACC依赖股权成本，而股权成本来自CAPM，CAPM自身存在大量问题；（4）忽略期权价值——灵活性、可延迟投资、可调整规模——这些"期权"价值DCF无法捕获；（5）杠杆公司的问题——高杠杆公司的WACC和FCF计算变得更复杂，不同假设下的结果差异巨大；（6）无形资产——品牌、专利、组织资本不体现在现金流量表上；（7）时间上的一致性——WACC和FCF都必须用名义值或实际值之一，混用会导致系统性偏差。\\
\textbf{一个有趣的替代方案：}Aswath Damodaran（纽约大学，被公认为"估值之父"）的解决方案是——把每个假设明确写出，定期更新，不做"点估计"而做"区间估计"。
}

\newpage
% ============================================================
\chapter{现代投资组合理论——Markowitz的均值—方差框架}
% ============================================================

\begin{scenario}
\textbf{一只股票 vs 十五只股票}

如果你在2019年初将全部资金买入某一只A股——比如中国平安——持有到2024年底的总收益约为-5\%。但如果你买的是沪深300指数（同时持有300只股票），同期的收益约为+15\%。如果你在300只中再分散到债券和黄金——收益更高、波动更小。分散化的好处不需要复杂的金融理论来理解——它是"不把所有鸡蛋放一个篮子里"的数学证明。
\end{scenario}

\begin{qualitative}
\textbf{分散化的本质：协方差}

分散化之所以能降低风险，不是因为股票数量多，而是因为不同资产的价格不会完全同步运动。当A跌时B可能涨（或跌得少）——这就是协方差和相关系数的力量。

\textbf{系统性风险 vs 非系统性风险：}
\begin{itemize}
  \item 非系统性风险（公司特有、行业特有）——可以通过分散化消除
  \item 系统性风险（市场整体、宏观经济、利率、战争）——无法通过分散化消除
\end{itemize}
\end{qualitative}

\begin{quantitative}
\textbf{双资产组合的预期收益和方差}

\[
E(R_p) = w_A E(R_A) + w_B E(R_B)
\]
\[
\sigma_p^2 = w_A^2 \sigma_A^2 + w_B^2 \sigma_B^2 + 2 w_A w_B \rho_{AB} \sigma_A \sigma_B
\]

当$\rho=1$时完全无分散化效果；当$\rho=-1$时理论上可消除所有风险（现实中几乎不存在完全负相关的资产）。

\textbf{有效前沿（Efficient Frontier）：}给定一组风险资产，所有达到"给定风险下的最高收益"或"给定收益下的最低风险"的投资组合构成有效前沿。有效前沿上的每个点对应一个最优组合。

\textbf{最小方差组合（Minimum Variance Portfolio）：}有效前沿上风险最低的点。

\textbf{夏普比率（Sharpe Ratio）：}
\[
SR = \frac{E(R_p) - R_f}{\sigma_p}
\]
最大夏普比率组合——引入无风险资产后的"最优"风险组合。
\end{quantitative}

\begin{lessonsummary}
\begin{enumerate}
  \item 分散化的力量来自资产之间的不完全相关性——不是来自数量本身
  \item 有效前沿展示了风险与收益之间的"最佳妥协"——任何一个有效前沿上的点都优于不在前沿上的点
  \item 15-20只充分分散的股票可以消除约90\%的非系统性风险
  \item 夏普比率是衡量风险调整后收益的标准化指标——投资的目标不是最大化收益，而是最大化单位风险的回报
\end{enumerate}
\end{lessonsummary}

\begin{discuss}
\begin{enumerate}
  \item 2008年金融危机期间，几乎所有资产价格同时暴跌——股票、债券、房地产、大宗商品（除了黄金）。当相关系数全都趋近于1时，分散化失效了。这是否暴露了Markowitz框架的致命缺陷？"危机时期的相关性突变"（Regime-Switching Correlation）如何应对？
  \item 在中国A股市场，投资者能否通过仅持有股票（不持有债券、商品等）实现有效分散化？A股内部的行业间相关系数大约是多少？（答案：约0.5-0.7，远高于美股约0.3-0.5）这是否意味着A股投资者需要比美股更"跨资产"的分散化？
  \item Markowitz获得了诺贝尔经济学奖（1990年），但他的博士论文答辩时Milton Friedman说"这不是经济学"——62年后他的框架仍然是资产管理行业的理论基础。一个"被当时最伟大的经济学家否定"的理论，如何成为了行业标准？
\end{enumerate}
\end{discuss}

\begin{exercises}
\begin{enumerate}
  \item 资产A：$E(R)=12\%$，$\sigma=20\%$；资产B：$E(R)=8\%$，$\sigma=10\%$；$\rho=0.3$。计算$w_A$从0到1变化时（步长0.1）各组合的$\sigma_p$和$E(R_p)$。描出有效前沿。最小方差组合的$w_A$是多少？
  \item 在（1）的基础上引入无风险资产$R_f=3\%$，计算最大夏普比率组合。
  \item 用Python从万得或Tushare下载5只A股（分属不同行业）近3年的日收益率数据，计算协方差矩阵，用数值方法求解有效前沿。（提示：使用scipy.optimize.minimize）
\end{enumerate}
\end{exercises}

\xuehaiinline{
\textbf{有效前沿的"输入敏感性问题"：}Michaud（1989）的经典论文《The Markowitz Optimization Enigma》指出——均值-方差优化对输入参数（预期收益）极度敏感。预期收益的微小变化会导致最优组合权重的大幅变化。这意味着：历史上（样本内）的"最优组合"到了未来（样本外）往往表现不佳。一个被称为"误差最大化器"的优化器——它倾向于给那些历史上表现好（但未来未必持续）的资产最大的权重。\\
\textbf{实践中怎么办？}（1）Black-Litterman模型——把市场均衡收益作为先验，加入投资者观点作为后验；（2）风险平价（Risk Parity）——不预测预期收益，只基于风险来分配权重；（3）等权重组合——简单粗暴但样本外表现经常优于优化组合（"1/N"之谜，DeMiguel et al. 2009）。\\
\textbf{相关系数的"体制转换"：}Longin \& Solnik（2001）发现国际股票市场之间的相关系数在熊市中显著高于牛市中。这意味着——正是你最需要分散化的时候（全球暴跌），分散化的效果最差。这个发现的含义：基于历史相关系数的分散化策略，在危机时提供的保护比预期更少。
}

\newpage
% ============================================================
\chapter{资本资产定价模型与套利定价理论}
% ============================================================

\begin{scenario}
\textbf{为什么高Beta股票在熊市中跌得最惨？}

2022年A股市场整体下跌约15\%。但不同股票的跌幅差异巨大——金融股（$\beta \approx 0.7$）平均跌约10\%，而科技股（$\beta \approx 1.5$）平均跌约25\%。市场"惩罚"高Beta股票在下跌市中的表现，但也"奖励"它们在牛市中更大的涨幅。CAPM解释了为什么：承担更多系统性风险的人，应该获得更高的预期回报——但这种"公平"只在统计意义上的长期成立，在任意给定的短期并不保证。
\end{scenario}

\begin{qualitative}
\textbf{CAPM的核心思想}

CAPM回答了一个问题：如果所有投资者都按Markowitz框架做分散化，那么一项资产的预期收益应该由什么决定？

\textbf{答案：}只有系统性风险（市场风险）被定价。非系统性风险可以被免费分散掉，所以市场不给它额外的回报。

\[
E(R_i) = R_f + \beta_i [E(R_m) - R_f]
\]

$\beta$衡量的是股票收益率对市场收益率变动的敏感度。$\beta=1$——跟市场同步；$\beta>1$——比市场更"激进"；$\beta<1$——比市场更"防御"。
\end{qualitative}

\begin{quantitative}
\textbf{$\beta$的估计}
\[
\beta_i = \frac{\text{Cov}(R_i, R_m)}{\sigma_m^2}
\]

通常用过去3-5年的月度或日收益率做回归。但历史$\beta$不一定是未来$\beta$——Blenman（2016）发现中国A股$\beta$的时间序列不稳定性远高于美股。

\textbf{证券市场线（SML）：}以$\beta$为横轴、$E(R)$为纵轴的直线。所有资产都应当落在这条直线上。落在线上方——$\alpha>0$（被低估），下方——$\alpha<0$（被高估）。

\textbf{詹森$\alpha$（Jensen's Alpha）：}
\[
\alpha_i = (R_i - R_f) - \beta_i (R_m - R_f)
\]

\textbf{Fama-French三因子模型：}
\[
E(R_i) - R_f = \beta_i [E(R_m) - R_f] + s_i \cdot SMB + h_i \cdot HML
\]
市场因子、规模因子（小盘股溢价）、价值因子（低P/B溢价）。

\textbf{Fama-French五因子模型}（2015）增加盈利因子（RMW）和投资因子（CMA）。
\end{quantitative}

\begin{lessonsummary}
\begin{enumerate}
  \item CAPM说——资产的预期收益只取决于其系统性风险（$\beta$），非系统性风险可以通过分散化免费消除
  \item CAPM的假设在现实中没有一条完全成立——但它作为一个"基准"仍然非常有用
  \item Fama-French多因子模型在解释截面收益率差异方面远优于CAPM
  \item 在中国A股市场，规模因子和动量因子的表现与美股有显著差异
\end{enumerate}
\end{lessummary}

\begin{discuss}
\begin{enumerate}
  \item Roll的批评（1977）：CAPM要求的"市场组合"应该是所有资产（包括人力资本、房地产、艺术品……）的加权组合，而不仅仅是股票市场指数。我们实际使用的"市场收益率"是次优的——这个批评是否动摇了CAPM的实用价值？还是说——所有模型都是错的，有些是有用的（Box, 1976）？
  \item 在中国A股市场，"小盘股溢价"（小公司长期跑赢大公司）是否存在？研究结论：1990-2015年存在，但2017年后消失了（甚至反转）——为什么？是因为注册制改革改变了壳价值、还是外资流入改变了定价逻辑？
  \item 在中国市场，"低波动率异象"（低$\beta$股票长期有更好的风险调整后回报）是否成立？如果成立，这对CAPM是一个更根本的挑战——因为CAPM预测高$\beta$应该有高回报，而实际恰恰相反。
\end{enumerate}
\end{discuss}

\begin{exercises}
\begin{enumerate}
  \item 无风险利率2.5\%，市场预期收益10\%，市场风险溢价7.5\%。三只股票的$\beta$：A（0.7）、B（1.0）、C（1.8）。按CAPM，各自的预期收益是多少？C的预期收益比A高多少？这个"额外收益"是否足以补偿其额外风险？
  \item 用过去3年的日收益率数据计算：中国平安(601318)、贵州茅台(600519)、宁德时代(300750)相对于沪深300指数的$\beta$。哪只最"防御"？哪只最"激进"？在2024年的行情中，CAPM的预测与实际表现匹配吗？
  \item 任选10只A股，计算每只的$\beta$和过去3年的年化收益率。用横截面数据检验：高$\beta$的股票是否有更高的平均收益？用散点图和回归线展示结果。
\end{enumerate}
\end{exercises}

\xuehaiinline{
\textbf{Fama vs. Campbell——两大学派的因子之争：}Eugene Fama（芝加哥大学，有效市场假说的旗手）认为多因子模型只是对风险的更准确度量——小盘股溢价是对"小公司风险"的补偿。而John Campbell（哈佛大学）认为部分因子实际上是行为偏误的结果——投资者系统性地高估了成长股的增长、低估了价值股的复苏，导致价值溢价。这两种解释有根本性的分歧——因子溢价来自风险补偿还是市场错误定价？如果是风险补偿，那么因子策略永远有效；如果是错误定价，那么随着更多投资者使用因子策略，"套利"会消除这些溢价。\\
\textbf{A股市场的特殊因子：}中国市场的因子研究中发现了一些"中国特色"因子——如"国有股占比"（国企vs民企的收益差异）、"政治关联度"（与政策方向的敏感度）、"散户关注度"（基于搜索量/论坛发帖量）。Liu, Stambaugh \& Yuan（2019）提出了一个"中国版三因子模型"——市场、规模、价值（用EP比率而非账面市值比），在中国A股的表现优于Fama-French原版。
}

\newpage
% ============================================================
\chapter{有效市场假说与行为金融}
% ============================================================

\begin{scenario}
\textbf{黑色星期一——一天22.6\%的暴跌没有任何新闻}

1987年10月19日，道琼斯工业平均指数一天暴跌22.6\%——这是史上最大的单日百分比跌幅，超过了1929年大萧条的任何一天。但当天并没有爆发战争、没有央行意外加息、没有恐怖袭击——没有任何"应该"导致市场暴跌20\%+的新闻。如果市场是有效的，这么大的价格变动必须对应同等量级的新信息——但这个信息并不存在。这个谜题至今没有一个被广泛接受的答案，但它对"市场总是理性的"这一假设提出了最尖锐的挑战。
\end{scenario}

\begin{qualitative}
\textbf{有效市场假说（EMH）的三种形式}

\textbf{弱式有效：}股价已反映所有历史价格和交易量信息。技术分析（K线、指标、形态）不能持续获得超额收益。

\textbf{半强式有效：}股价已反映所有公开信息（财报、新闻、公告）。基本面分析不能持续获得超额收益。

\textbf{强式有效：}股价已反映所有信息——包括内幕信息。任何人，包括公司内部人，都不能持续获得超额收益。

实证检验的结论：弱式有效在大多数市场基本成立（至少趋势跟踪策略的收益不足以覆盖交易成本）；半强式有效受到大量"异象"的挑战（盈余公告后漂移、动量效应等）；强式有效被内幕交易查处案例系统性否定。
\end{qualitative}

\begin{quantitative}
\textbf{事件研究法（Event Study）}

用于检验半强式有效的标准方法：
\begin{enumerate}
  \item 定义事件窗口（如盈余公告日前后的[-5, +5]交易日）
  \item 计算累计异常收益（CAR）——个股收益率减去预期收益率（通过CAPM或市场模型估算）
  \item 检验CAR是否在统计上显著异于零
\end{enumerate}

\textbf{盈余公告后漂移（PEAD）：}Ball \& Brown（1968）发现——盈余超预期的股票在公告日后数周内继续跑赢市场，盈余低于预期的股票继续跑输。这直接违背了半强式有效——如果股价已经反映了公开信息，为什么信息公布后还会继续"漂移"？
\end{quantitative}

\begin{lessonsummary}
\begin{enumerate}
  \item 市场有效性的三种形式——弱式、半强式、强式——代表了市场效率的不同层次
  \item EMH是一个"基准"而非"真理"——它描述了如果市场完全理性会是什么样子，但现实中市场并不完全理性
  \item 行为金融学发现了大量EMH无法解释的异象——但大部分异象在扣除交易成本和风险调整后消失或大幅减弱
  \item Fama和Shiller分享了同一个诺贝尔奖（2013年）——一个相信市场永远有效，一个相信市场经常非理性
\end{enumerate}
\end{lessonsummary}

\begin{discuss}
\begin{enumerate}
  \item 在一个不完全有效的市场中，主动投资（自己选股、选基金）的合理性是什么？如果90\%的主动基金经理在10年维度跑不赢指数，那么主动投资是一个"零和博弈"甚至"负和博弈"（扣除管理费后）。为什么还有上万亿资金在主动管理中？
  \item "幸存者偏差"的问题——我们看到的"股神"（Buffett, Lynch, Simons）是因为他们真的有能力，还是因为失败的投资者已经从公众视野中消失了（Survivorship Bias）？如何区分"技能"和"运气"？
  \item 行为金融偏误——"过度自信偏误"让个人投资者频繁交易，研究显示（Barber-Odean, 2000）频繁交易的投资者年化收益比市场平均低约6-7\%。为什么过度自信在男性投资者中比女性投资者更显著？
\end{enumerate}
\end{discuss}

\begin{exercises}
\begin{enumerate}
  \item 选取过去3年A股市场的5次重大政策公告（如降准、降息、印花税调整），用事件研究法分析每次公告前后[-5, +10]个交易日的市场反应模式——价格是瞬间到位还是持续调整？
  \item 设计一个简单的动量策略——每月买入过去6个月涨幅最高的前10\%股票，卖空跌幅最高的前10\%股票。回测过去5年的收益率。扣除交易成本（千分之三）后，这个策略还能盈利吗？Carhart（1997）的动量因子在A股市场有效吗？
  \item 在线搜索"前景理论"的实验——复制Kahneman-Tversky的经典损失厌恶实验，在班级中进行调查，统计数据支持"损失厌恶系数约2"的结论吗？
\end{enumerate}
\end{exercises}

\xuehaiinline{
\textbf{Kahneman的快慢思维与投资决策：}Daniel Kahneman在其著作《思考，快与慢》中将人类思维分为两个系统——系统1（快速、直觉、情绪化）和系统2（缓慢、理性、分析性）。大多数投资决策是由系统1做出的——因为系统2太消耗精力。当市场剧烈波动时，系统1接管——恐慌性卖出或贪婪性买入。\\
\textbf{处置效应（Disposition Effect）：}Shefrin \& Statman（1985）发现投资者倾向于过早卖出盈利的股票（锁定收益——系统1的风险厌恶）和过久持有亏损的股票（期待回本——系统1的损失厌恶之后的"风险偏好"）。Odean（1998）使用美国某大型折扣券商的交易数据验证了这个效应。中国A股市场同样存在处置效应——且在散户中比机构中更显著。\\
\textbf{Richard Thaler的"心理账户"：}为什么人们对"工资"和"奖金"采取不同的消费倾向？——心理账户把不同来源的钱分开对待。同一笔1000元——如果是工资，可能存起来；如果是意外之财（红包、彩票），会立即花掉。从经济学角度，钱是可替代的——1000元就是1000元，不管它来源是什么。行为上，人不是完全理性的"经济人"。
}

\newpage
% ============================================================
\chapter{中央银行与货币政策}
% ============================================================

\begin{scenario}
\textbf{一场"以利率为武器"的全球实验}

2022-2023年，美联储在一年内加息11次——从0.25\%加到5.5\%，幅度为40年来最大。但同时期中国央行却在降息。两国面临截然不同的经济挑战：美国通胀（CPI）一度超过9\%（40年最高）；中国的是内需不足、房地产低迷、甚至出现了通缩压力。同一时间、同一全球经济环境下，两大央行的决策方向完全相反——谁的决策是对的？政策的传导效果如何评估？
\end{scenario}

\begin{qualitative}
\textbf{央行的三个使命}

现代央行通常承担三个相互冲突的使命：
\begin{enumerate}
  \item 价格稳定（控制通胀/通缩）
  \item 充分就业（促进经济增长）
  \item 金融稳定（防止系统性危机）
\end{enumerate}

三个目标之间的冲突在2022年暴露无遗——为了降通胀（目标1）必须加息，但加息可能引发经济衰退（损害目标2）和金融风险（损害目标3）。

\textbf{货币政策的传导渠道：}
\begin{itemize}
  \item 利率渠道——加息→借钱变贵→投资消费下降→通胀下行
  \item 信贷渠道——加息→银行收缩信贷→企业融资困难→经济放缓
  \item 资产价格渠道——加息→股价/房价下跌→财富缩水→消费减少
  \item 汇率渠道——加息→本币升值→出口竞争力下降
\end{itemize}
\end{qualitative}

\begin{quantitative}
\textbf{Taylor Rule（泰勒规则）}

估算政策利率的"合理水平"：
\[
i = r^* + \pi + 0.5(\pi - \pi^*) + 0.5(y - y^*)
\]

其中$r^*$为中性实际利率，$\pi$为当前通胀率，$\pi^*$为目标通胀率，$y-y^*$为产出缺口（实际GDP与潜在GDP的偏离百分比）。

\textbf{数量方程式——货币数量与通胀的关系：}
\[
MV = PY
\]

$M$=货币供应量，$V$=货币流通速度，$P$=价格水平，$Y$=实际产出。这个公式是货币主义（Milton Friedman）的理论基础——"通胀在任何时候、任何地方都是一种货币现象"。

\textbf{货币乘数：}
\[
m = \frac{1}{rr}
\]

\textbf{中国的实际情况：}中国的货币政策传导效率低于美国——原因：利率市场化尚未完全完成、部分借款主体（国企、地方政府平台）对利率不敏感、商业银行的风险偏好受监管约束。2023年中国货币乘数约7.5（准备金率约7.4\%），但实际M2增速（约9-10\%）主要由信贷扩张驱动而非乘数效应。
\end{quantitative}

\begin{lessonsummary}
\begin{enumerate}
  \item 央行的三个使命相互冲突——价格稳定、充分就业、金融稳定不可能同时完美实现
  \item Taylor Rule提供了政策利率的量化基准，但其参数选择存在高度主观性
  \item 货币政策的传导效率取决于一国的金融体系结构——中国的信贷渠道强于利率渠道
  \item 现代货币理论（MMT）对传统央行的观点提出了根本性挑战——但其在通胀环境下的适用性存疑
\end{enumerate}
\end{lessonsummary}

\begin{discuss}
\begin{enumerate}
  \item 2010-2020年，日本央行的资产负债表从GDP的20\%扩张到130\%以上——但通胀从未超过2\%。这是否意味着"货币数量论"在现代金融体系中已经失效？还是说货币流通速度（V）的大幅下降抵消了基础货币（M）的扩张？
  \item 美联储的"平均通胀目标制"（AIT, 2020年引入）——允许通胀在一段时间内适度超过2\%来弥补之前低于2\%的时期。这个政策在2021-2022年通胀飙升后受到了严厉批评——它是刺激了通胀预期失控的罪魁祸首吗？
  \item 央行的独立性是否应该被保障？——如果政府可以命令央行"印钞发债"，短期刺激经济但长期导致恶性通胀（如委内瑞拉、津巴布韦）。独立的央行更关注长期稳定而不是短期政治诉求——但它也更"官僚化"和"反应迟缓"？
\end{enumerate}
\end{discuss}

\begin{exercises}
\begin{enumerate}
  \item 中国当前CPI约0.5\%（2024年），潜在GDP增速约5\%，产出缺口约-2\%。目标通胀3\%（中国央行长期目标），中性实际利率$r^*$约2\%。用Taylor Rule，当前政策利率的"合理水平"应该是多少？与实际的1年期LPR（约3.1\%）对比，结论是什么？
  \item 计算M2增速与名义GDP增速的差值——如果M2增速（约9\%）持续高于名义GDP增速（约5\%），长期来看对通胀意味着什么？为什么中国过去10年M2/GDP从180\%升到220\%以上但通胀并没有失控？
  \item 选取近3次中国央行降准事件，分析每次降准后1个月内股市、债市和房地产市场的反应模式。降准是"利好"还是"预期内的利多出尽"？
\end{enumerate}
\end{exercises}

\xuehaiinline{
\textbf{现代货币理论（MMT）的争议：}MMT的核心主张——主权货币政府（如美国、日本、中国）不存在"破产"的可能，因为它可以无限创造自己的货币来偿还以本币计价的债务。因此，政府债务的上限不是"能不能还"，而是"会不会引发通胀"。这个理论被左翼政治势力热烈拥抱——它意味着政府可以无限增加支出用于公共投资、全民医保、绿色转型而不必担心"没钱"。\\
\textbf{批评者（包括大多数主流经济学家）的反驳：}（1）MMT在通胀压力下的处方——增税来吸收多余的购买力——在政治上比"印钞"更难实施；（2）日本的经验不可复制——日元的国际储备地位、长期通缩心态、深度老龄化的特殊背景使日本成为"特例"而非"通则"；（3）当一个国家试图"印钞"来为财政赤字融资时，国际投资者可能会抛售该国国债（如英国2022年的"迷你预算危机"——特拉斯减税方案引发英镑暴跌、国债收益率飙升）。\\
\textbf{中国的实践：}中国确实在2020年发行了1万亿抗疫特别国债，央行并未直接"购买"（中国《中国人民银行法》禁止央行直接为财政融资），但通过降准和MLF投放了大量流动性。中国式的"准MMT"——通过国有银行体系间接承担了类似MMT的功能——但体制性地避免了恶行通胀，原因在于中国银行信贷的流向受到强有力调控。
}

\end{document}
